\documentclass{article}
\usepackage[left=2cm, right=2cm, top=2cm, bottom=2cm]{geometry}
\usepackage{graphicx}
\usepackage{amsmath}
\usepackage{amssymb}
\usepackage{amsthm}
\usepackage{fancyhdr}
\usepackage{verbatim}
\usepackage{listings}
\usepackage{xcolor}
\usepackage{pdfpages}
\usepackage{cancel}
\usepackage{physics}
\usepackage{esint}

\lstset{
    language=C++,
    basicstyle=\ttfamily\footnotesize,
    keywordstyle=\color{blue},
    commentstyle=\color{green},
    stringstyle=\color{red},
    numbers=left,
    numberstyle=\tiny\color{gray},
    frame=single,
    breaklines=true,
    captionpos=b,
}


\title{MTH 553: Lecture \# 36}
\author{Cliff Sun}

\newtheorem{theorem}{Theorem}[section]
\newtheorem{lemma}[theorem]{Lemma}
\newtheorem{definition}[theorem]{Definition}
\newtheorem{conjecture}[theorem]{Conjecture}
\newtheorem{proposition}[theorem]{Proposition}
\newtheorem{corollary}[theorem]{Corollary}
\newtheorem{one minute paper}[theorem]{One Minute Paper}

\pagestyle{fancy}
\lhead{\textbf{\thepage}\ \ \nouppercase{\rightmark}}
\chead{MTH 553: Lecture \# 36}
\rhead{Cliff Sun}

\begin{document}

\maketitle

\section*{Lecture Span}
\begin{enumerate}
    \item Dirichlet's Principle
\end{enumerate}

Here, we attempt to solve Poisson's equations by minimizing energy. Let $\Omega$ be bounded and smooth in $\mathbb{R}^n$. Let $f \in C\left( \overline{\Omega} \right)$ and $g \in C(\partial \Omega)$. 

Define 

\begin{equation}
    E[w] = \int_{\Omega}\left( \frac{1}{2}|\nabla w|^2 - w f \right)dx
\end{equation}

For $w$ in the set 

\begin{equation}
    A = \{w \in C^2(\overline{\Omega}): \quad w = g \text{ on $\partial \Omega$}\}
\end{equation}

Here, we assume that $A \neq \emptyset$. Note that $w$ is really any arbitrary function that equals $g$ on $\partial \Omega$. 

\begin{theorem}
    If 
    \begin{align}
        -\triangle u &= f, \quad \Omega \\
        u &= g, \quad \partial \Omega
    \end{align}
    Then, Dirichlet's principle states that:
    \begin{equation}
        E[u] = \min_{w \in A}E[w]
    \end{equation}
\end{theorem}

Note that if $u \in A$ and also satisfies Dirichlet's Principle, then 

\begin{align}
        -\triangle u &= f, \quad \Omega \\
        u &= g, \quad \partial \Omega
\end{align}

Therefore, 

\begin{center}
    DP $\iff$ $u$ solves Poisson's equations
\end{center}

\begin{proof}
    (a) Let $-\triangle u = f$. Then for all $w \in a$, then we have 
    \begin{align}
        0 &= \int_{\Omega}\left( -\triangle u - f \right)\left( u - w \right)dx \\
        &= \int_{\Omega}\left[ \nabla u \cdot \nabla (u - w) - f(u - w) \right]
    \end{align}
    By green 1 since $ u = g = w$ on $\partial \Omega$. Then 
    \begin{align}
        &= \int_{\Omega}\left[ |\nabla u|^2 - \nabla u \cdot \nabla w - fu + fw \right]dx \\
        &\geq \int_{\Omega}\left[ \frac{1}{2}|\nabla u|^2 - \frac{1}{2}|\nabla w|^2 - fu + fw \right]dx
    \end{align}
    Therefore, 
    \begin{equation}
        E[w] \geq E[u]
    \end{equation}
    Which implies the Dirichlet principle. 
\end{proof}

We next prove $\impliedby$. 

\begin{proof}
    (b) Fix $v \in C^{\infty}_0(\Omega)$ and let 
    \begin{equation}
        e(\tau) = E[u + \tau v], \quad \tau \in \mathbb{R}
    \end{equation} 
    Since $u + \tau v \in A$, the DP implies 
    \begin{equation}
        e(0) \leq e(\tau) \implies e'(0) = 0
    \end{equation}
    Then 
    \begin{align}
        0 &= \frac{d}{dt}\int_{\Omega}\left[ \frac{1}{2}|\nabla u|^2 + \taiu \nabla u \cdot \nabla v + \frac{\tau^2}{2}|\nabla v|^2 - uf - \tau v f\right] \\
        &= \int_{\Omega}\left[ \nabla u \cdot \nabla v - vf \right] \\
        &= \int_{\Omega}\left( -\triangle u - f \right)v dx
    \end{align}
    By green 1 since $v = 0$ on $\partial \Omega$. Therefore, for all $v \in C^{\infty}_0(\Omega)$, we have that $-\triangle u - f = 0$ and $u \in A$, therefore, $u = g$ on $\partial \Omega$. 
\end{proof}

\end{document}